\begin{table}[t]
\centering
\caption{Hard case analysis: Where ArchaeoGPT succeeds but ViT fails.}
\label{tab:hard_cases}
\small
\begin{tabular}{p{2cm} p{2.5cm} p{2.5cm} p{4.5cm}}
\toprule
True Culture & ArchaeoGPT & ViT & Description \\
\midrule
苏贝希文化晚段 & 苏贝希文化晚段 (0.83) & 穷科克文化早段 (0.39) & 侈口，口沿一侧较高，高领，颈肩单宽带耳，鼓腹，圜底，黄地红彩。口沿下至器底绘一周竖线纹，竖线纹两侧条带内分别绘相对的平行斜线纹，其中有两条条带内绘竖向交错的三角... \\
大汶口文化 & 大汶口文化 (0.98) & 仰韶文化早期 (0.57) & 泥质青灰陶。器体轮制，通体磨光。侈口，圆唇，折沿，深圆腹，圜底较缓，喇叭形圈足。口沿内侧一圈涂红彩，圈足饰一周凸棱。... \\
马家窑文化半山期 & 马家窑文化半山期 (0.90) & 马家窑文化马厂期 (0.50) & 泥质土黄陶。侈口，圆唇，短颈，鼓腹，小平底。双腹耳。施红、黑彩。口沿内绘垂弧纹、锯齿纹；外颈部绘一周锯齿纹；腹部绘锯齿菱形方格纹，内填十字圆点纹；下腹绘一周宽带... \\
侯家寨二期文化 & 侯家寨二期文化 (0.80) & 黄瓜山文化 (0.42) & 细泥黄陶。敛口，弧腹。器表有黑彩纹饰，口沿处饰一周条纹，肩部饰两周条纹，口沿下饰圆点纹，间以三道竖线。... \\
菜园文化 & 菜园文化 (0.73) & 大汶口文化 (0.45) & 泥质橙黄陶。侈口，圆唇，束颈，鼓腹，平底稍内凹。罐底有压印纹。口沿内侧饰红彩短条带纹，颈部及上腹饰有红彩网格纹，上腹饰网格纹间用三道竖线纹分隔。腹径最大处有一周... \\
\bottomrule
\end{tabular}
\end{table}